% Chapter 6 merges the former Discussion and Conclusions chapters, in the
% order the structure guide prescribes: Summary, Contributions, Limitations,
% Future Work, Conclusions.
This chapter summarizes what the thesis found, states what it contributes, sets
out the limitations that qualify those contributions, sketches the natural
extensions the work opens up, and closes with the conclusions.
Sections~\ref{sec:disc-methodological}--\ref{sec:disc-joint} are the summary
and discussion: the methodological study, the clinical study, and what the two
say together.

\paragraph{Recap of what is being discussed.}
Before turning to what the results mean, a brief reminder of what the thesis
contributed and where each contribution was tested. The methodological work
introduced \textbf{INTERVenE}, a transformer family that reads knowledge-based
temporal-abstraction (KBTA) intervals as its native token stream --- an
autoregressive decoder (INTERVenE-Ar) that generates a full abstraction
trajectory with per-step risk readouts, and a bidirectional encoder
(INTERVenE-Enc) that yields single-pass joint risk and time-to-event
prediction. Both variants share a hierarchical concept embedder, continuous
time encoding, and static-context conditioning, and both are probed with
gradient-decoupled sparse-transcoder attribution so explanations resolve to
named clinical concepts. The clinical work formalized fourteen recommendations
of the national inpatient diabetes guideline as quality-assessment temporal
patterns over the same interval stream, so measurement and prediction read one
representation. Three experiments carried the argument:
Chapter~\ref{chapter:results} reports (i) the MIMIC-IV benchmark that compares
INTERVenE against GRU-D, STraTS in two configurations, and KarmaLego+TPF on
$57{,}078$ diabetic admissions; (ii) the Kineret adherence analysis across
$48{,}354$ admissions in six Israeli general hospitals, before and after the
national implementation effort, with propensity-weighted outcome regressions;
and (iii) the Kineret seven-arm ladder that transfers INTERVenE-Enc to a
second cohort, isolates the knowledge-base lift against $\sigma$-bin
discretization, and tests whether the measured adherence signal helps
prediction. The discussion that follows reads those three results together.

% Sections 6.1-6.3 of Chapter 6. The chapter opening, and the Contributions,
% Limitations and Conclusions sections, live in content/summary/summary.tex.
\section{The Methodological Study}
\label{sec:disc-methodological}

\subsection{Where the Predictive Gain Comes From}

Two results have to be read together, and doing so gives a sharper conclusion than either alone.

INTERVenE-Enc outperforms every baseline on the shared benchmark, including KarmaLego+TPF, which consumes the identical interval stream. That comparison is the one that isolates modeling from representation, and it is the largest gap in the table. Yet replacing the KBTA vocabulary with a non-semantic $\sigma$-bin discretization leaves discrimination essentially unchanged.

The natural reading is that the predictive gain comes from the transformer operating over persistent intervals, not from the intervals carrying clinical names. Collapsing every interval to a one-second point event confirms the other half: performance falls sharply, so the architecture depends on duration and on state concurrency rather than acting as a point-event detector. What the model exploits is that a symbol \emph{holds} over a span and that spans \emph{overlap} --- structure a discretization supplies whether or not a clinician chose its boundaries.

This is a negative result for the claim a knowledge-based thesis might most want to make. Clinical knowledge does not buy predictive accuracy here. It is worth stating plainly rather than softening, because the alternative claim it leaves standing is the more useful one: the interpretability is obtained at essentially no accuracy cost. Interpretability bought by sacrificing performance is rarely adopted; interpretability that is free is a different proposition.

The representation helps most where the target is itself an abstraction. Severe dysglycemia outcomes are defined by threshold-crossing states and their durations, so the input injects the clinical definition directly. The gain also extends to death and to length of stay, which no single KBTA threshold defines --- there the model appears to combine abstracted states such as hypoalbuminemia, organ-stress trends, and acid--base shifts into a composite signal of deterioration, consistent with the attribution results.

\subsection{What the Explanations Do and Do Not Show}

Because every input token is a named clinical concept, attribution resolves to clinical vocabulary by construction rather than by subsequent interpretation. The death drivers are clinically coherent, and one detail is more informative than the overall coherence: the model assigns opposite signs to an elevated enzyme level and to a decreasing enzyme trend. A method keyed to how often a concept appears could not make that distinction, so the attribution is tracking clinical direction rather than clinical activity.

Two limits are structural and were carried into the reporting. Gradient attribution measures local sensitivity, not causal necessity. And it describes the fitted function, not the patient --- it identifies what drove a prediction, which is a claim about a model. Neither limit is specific to sparse transcoders; both apply to saliency, integrated gradients, and SHAP alike. What the KBTA vocabulary changes is not the epistemic status of the attribution but its legibility: the same claim, expressed in concepts a clinician can evaluate.

\subsection{Two Variants, Two Niches}

The variants are complementary rather than redundant. INTERVenE-Enc is the production predictor: strongest single-pass accuracy, and low enough latency to score a full test cohort in seconds. INTERVenE-Ar is the inspection tool: it generates the abstraction trajectory itself, so its explanation is its output, and only it localizes \emph{when} and after \emph{which} events risk rises.

The roughly $77\times$ latency difference follows directly from generation and is not an implementation shortcoming. It does mean the decoder's value is diagnostic rather than operational, and its weaker performance on death --- where a rollout must stay coherent over many steps before a sparse terminal token appears --- marks the limit of the generative readout rather than of the representation.

\section{The Clinical Study}
\label{sec:disc-clinical}

\subsection{Principal Findings}

Adherence improved after the national recommendations were introduced, and higher adherence was consistently associated with lower mortality and reduced hyperglycemic complications. But adherence was not uniform, and the pattern of non-uniformity is the finding.

The split follows the \emph{kind} of action a recommendation requires, not its clinical weight. Recommendations satisfied by a single discrete action --- ordering a test, discontinuing a medication on admission --- are near-universally implemented. Those requiring repeated insulin decisions across a stay are not, several with a median of zero, meaning the expected action was absent altogether in most eligible admissions. The improvement following implementation was significant but modest, its $P$ value reflecting the number of admissions rather than the size of the shift, and it came with substantial inter-hospital heterogeneity these data do not explain.

The clearest outcome signal is for $30$-day mortality. The propensity-weighted estimate is the primary one, being the only estimate whose exposure window is guaranteed to precede the outcome. Two descriptive estimates agree in direction --- whole-admission adherence adjusted for baseline characteristics, and a graded fall in mortality across the adherence range. A graded association is harder to explain by case mix alone, though it shares the timing concern noted below. Because the three use different exposure definitions, they corroborate direction rather than magnitude.

The glycemic findings show a favorable risk profile rather than an unambiguous benefit. Greater adherence bought a large reduction in hyperglycemic burden --- roughly ninefold larger on the log-odds scale than the accompanying rise in hypoglycemia --- at the cost of a small increase in mild hypoglycemia and no detectable increase in severe hypoglycemia. That is precisely the balance inpatient insulin recommendations are designed to strike. Since non-significance is not evidence of absence, the severe-hypoglycemia result shows no signal of harm at this sample size rather than establishing that none exists.

\subsection{Implications}

The most actionable pattern is a convergence. The recommendations with the largest implementation gaps --- insulin following hyperglycemia, and both components of routine basal dosing --- are also among those most consistently associated with outcomes after balancing measured baseline characteristics. The components where practice departs most from the guideline are therefore not peripheral ones.

This identifies them as natural targets for decision support, and for a specific reason: they are repeated decisions distributed across a hospitalization, which is exactly the kind of task where a prompt at the point of care can plausibly help, and unlike the admission recommendations they are nowhere near ceiling. A reminder to order creatinine would change almost nothing, because creatinine is already ordered.

This should not be read as an attribution of fault. A low adherence score records that the action the guideline expects is absent from the record --- not that the decision was wrong. A patient not eating, an imminent procedure, a deliberate decision to tolerate higher glucose in a frail patient: each is frequently the correct reason for a departure, and none is recoverable from records collected for clinical rather than research purposes. The measure is a population-level statistic, and so is the inter-hospital variation. What this work offers is a measurement instrument --- an automated, component-level, transferable way to ask which recommendations were followed and which track outcomes --- not a decision-support system.

Nothing restricts the approach to recommendations a guideline already contains. Because each is declared externally as a pattern rather than implemented in analysis code, a practice merely \emph{suspected} to matter costs no more to evaluate than the patterns scored here. That inverts the usual question: one could ask which candidate practices earn a place in a guideline, and --- more uncomfortably --- which current recommendations track no outcome at all. Such candidates would be hypotheses for prospective evaluation, and a wide search would need more conservative multiplicity correction than was applied here.

\section{What the Two Studies Say Together}
\label{sec:disc-joint}

The two studies share a representation rather than an experiment, and each supplies something the other cannot.

The methodological study establishes that a transformer over KBTA intervals is competitive with the established alternatives and that its explanations resolve to clinical concepts --- but on a public benchmark, with a knowledge base built for it, and with no notion of whether care followed a guideline. The clinical study establishes that guideline adherence can be measured automatically at national scale, that it improved unevenly, and that the least-implemented recommendations are among those most associated with outcomes --- but analyzes adherence statistically, as an exposure, without asking whether a model can use it.

What joins them is that adherence is computed over the same intervals the model reads. The Kineret experiment of Section~\ref{sec:evalm-kineret-task} closes that loop, and it is the one place where the thesis's two halves become a single question: does a measured quality-of-care signal carry predictive information beyond the clinical record it was derived from?

That experiment has now been run as a closing addition (Section~\ref{sec:eval-kineret-models}), with every arm at the MIMIC-IV configuration and no tuning. Two results follow.

The architecture transfers, and the knowledge base carries more predictive weight here than on MIMIC-IV. Substituting $\sigma$-bins costs $0.073$ AUPRC on Kineret but was within noise on MIMIC-IV. The knowledge base was authored on Kineret and transferred to MIMIC-IV, so the asymmetry is unsurprising: it is better matched to the cohort it was built for, where it also instantiates more fully --- eighteen of its concepts have no counterpart in the MIMIC-IV extract (Section~\ref{sec:method-preprocessing}). The MIMIC-IV finding therefore bounds what clinical knowledge contributes when ported to a new source, not what it contributes generally.

Neither cohort licenses the stronger claim that data-driven abstraction cannot work. The knowledge base is a hand-curated artifact and the $\sigma$-bins a fully automated one, so the contrast is between two abstraction methods at very different levels of effort, not between knowledge and data as such. A discretization tuned to the cohort, rather than applied off the shelf, was never tested.

The adherence signal does not improve the aggregate metrics, which was the likely outcome given that adherence is computed from the intervals the model already reads. Attribution shows the model does use it: compliance patterns rank among the top drivers for seven of eight outcomes. The signal is visible to the model but redundant with what it already has --- a statement about representation, not about whether following the guideline matters clinically, which Section~\ref{sec:eval-clinical} addresses separately.

A single caution applies to whatever those results show. Adherence is derived from the same record the model reads, so a predictive lift would demonstrate that the aggregation makes an existing signal easier to use, not that new information was introduced. That is a real and useful finding --- feature construction routinely earns its keep that way --- but it is a claim about representation, not about the clinical value of following the guideline. The latter question is what the propensity-weighted analysis addresses, and the two should not be conflated.


\section{Contributions}
\label{sec:contributions}

\subsection{To Computer Science}

\textbf{A transformer architecture over temporal abstraction intervals.}
INTERVenE takes symbolic intervals as its native input rather than raw
measurements: each interval becomes a START and an END token embedded through a
four-level concept hierarchy, with elapsed time entering through Time2Vec and
temporal rotary embeddings and static context through AdaLN-Zero. Two variants
were developed and compared --- a bidirectional encoder for single-pass risk and
time-to-event prediction, and an autoregressive decoder that generates
abstraction trajectories with per-step risk curves. The encoder outperforms the
established neural and symbolic alternatives on a shared benchmark, and leads
every comparator on both cohorts it was evaluated on.

\textbf{Evidence about what interval representations contribute.} Holding the
architecture fixed and varying only the input isolates two separable effects:
collapsing intervals to point events degrades performance sharply, while
replacing clinically named intervals with an automatically derived $\sigma$-bin
discretization does not. On MIMIC-IV the gain therefore belongs to persistence
and concurrency rather than to clinical naming --- a result that bounds what a
knowledge base buys a predictor, and that makes the architecture applicable to
interval-structured data outside medicine. The Kineret repetition qualifies it:
there the substitution cost accuracy, so the equivalence is demonstrated for a
source the knowledge base was ported to, not for one it was authored on.

\textbf{Attribution that resolves to named concepts.} Sparse transcoders attach
through a gradient-decoupled hook whose injected term is numerically zero, so
the forward pass is bit-exact and the logits explained are the deployed model's
own. Because the vocabulary consists of named concepts, the explanations are
clinical statements rather than feature indices.

\textbf{A released software stack.} A reimplementation of the Temporal Mediator
that performs abstraction and adherence scoring in one pass over a TAK
repository; a KarmaLego interval-mining pipeline with a Temporal Probabilistic
Profile representation; and INTERVenE with its attribution pipeline. All are
pure Python and deployable in ordinary environments, including the regulated
hospital platform this work ran inside.

\subsection{To Medicine and Medical Informatics}

\textbf{A single-schema knowledge base for inpatient diabetes care.} The $174$
TAK units encode the clinical semantics of inpatient care and the fourteen
recommendations of the Israeli national position paper as quality-assessment
patterns. Declarative and quality-assessment knowledge live in one
representation executed by one engine in a single pass, rather than being split
across a declarative formalism and a procedural one. Because it names clinical
concepts rather than dataset columns, it transfers to any EHR within the same
clinical scope once a concept index is supplied.

\textbf{An automated, component-level measure of guideline adherence.}
Compliance becomes a continuous quantity per recommendation and per admission,
comparable across hospitals and over time without chart review. Applied at
national scale across six hospitals, it shows that adherence levels differ by
the kind of action a recommendation requires rather than by its clinical weight.

\textbf{Evidence on which recommendations matter.} Under propensity weighting,
above-median adherence carries substantially lower odds of $30$-day mortality
and of hyperglycemia against a small rise in hypoglycemia. The recommendations
that were least adhered to are among those most strongly associated with
outcomes --- the thesis's most actionable finding, and a direct target for
decision support.

\textbf{Guideline adherence treated as an explanatory variable.} Procedural
quality-of-care knowledge is brought to bear on patient outcomes alongside the
declarative knowledge such analyses normally rely on, both statistically as an
exposure and as model input.

\section{Limitations}
\label{sec:limitations}

\subsection{The Representation}

The central limitation is intrinsic to the representation. The KBTA layer is only as expressive as the knowledge base behind it: a concept with no abstraction rule is simply absent from the model input. Extending INTERVenE to a new disease, population, or care setting therefore requires authoring a new knowledge base rather than retraining on new labels --- an expert-driven process, and substantially more laborious than adapting an end-to-end model.

A related limit is granularity. Rule-based abstraction deliberately collapses clinically equivalent values, so glucose of $250$ and of $600$\,mg/dL both map to severe hyperglycemia. The $\sigma$-bin ablation shows this costs little predictively, which is a reassurance about the architecture rather than about the knowledge base.

Beyond the representation, the methodological study evaluates one MIMIC-IV cohort. External validation, subgroup fairness across demographics and care settings, and attribution stability across random initializations all remain open. Reported intervals are patient-level bootstraps, not multi-seed variance: full training takes roughly $8$--$12$ hours for the encoder and over $24$ for the decoder on a single device, which put multi-seed sweeps out of reach. The intervals therefore describe sampling uncertainty for the model that was trained, not the variability of the procedure that produced it.

\subsection{The Clinical Study}

The design is observational and none of the estimates is a causal effect. Weighting balances measured covariates only; unmeasured factors cannot be excluded, and two matter especially --- real-time illness severity, and decisions to limit or de-escalate treatment.

The mortality analyses divide by exposure definition, and the division is worth being explicit about. The component-level and propensity-weighted models score adherence over a fixed early-admission window, so exposure strictly precedes outcome. The adjusted estimate instead averages adherence over the whole admission and is descriptive: because $30$-day mortality is predominantly an end-of-admission event, and guideline-concordant management is frequently de-escalated in patients on a deteriorating trajectory, that average incorporates the period immediately preceding death, where adherence may reflect the course of care rather than determine it. Dichotomizing at the median, as weighting requires, additionally discards dose-response information and costs power.

Adherence and outcomes were ascertained during the same hospitalization, and detecting a glycemic event depends on glucose being measured. Routine monitoring was therefore excluded from the glycemic models, but residual detection effects may remain for other monitoring recommendations --- visibly so for admission monitoring, which appears associated with \emph{higher} recorded hyperglycemia and is reported as an ascertainment artifact rather than as harm.

The pre--post comparison is uncontrolled, so secular trends in inpatient care cannot be separated from the implementation itself. The non-mortality models do not account for the competing risk of death, which biases those endpoints conservatively. The findings come from adult inpatients in Israeli general hospitals under a single national guideline, so both the adherence levels and the components that proved most consequential may differ elsewhere. Finally, the data are assembled from admissions rather than from a longitudinal record: ascertaining eligibility and glycemic outcomes from measurements rather than diagnosis codes limits the damage, but the comorbidity score and the proportion of hyperglycemic admissions carrying a documented diagnosis are both likely understated.

\subsection{Limits Common to Both}

Two further limits cut across the thesis. The knowledge base is portable in
principle because it names clinical concepts rather than dataset columns, but
landing it on a new source still costs a concept index: Kineret carried no
mapping to an international coding system at the time of this work, so roughly
$1{,}500$ local codes were resolved onto the vocabulary by hand. That labor is
site-specific and must be repeated at the next network, and where local codes
could not be identified it degraded the representation itself --- the
comorbidity burden entering the static vector is a \emph{modified} Charlson
index because the source supports nothing stricter.

And adherence measures the record, not the decision. Bedside reasoning is
largely invisible to the measure and always will be, since it is not the kind of
thing an electronic record captures. Adherence is also expected to decline as a
patient approaches death, when guideline-concordant management is deliberately
de-escalated, which is why the mortality analyses that carry weight score
adherence over a fixed early window rather than over the whole admission.

\section{Future Work}
\label{sec:future-work}

Several directions follow from the thesis but were not pursued in it. Some
address the limitations named above directly; others are extensions the work
opens up rather than repairs.

\textbf{Weighted aggregation up the guideline hierarchy.} Recommendations are
currently averaged with equal weight, so a critical requirement counts the
same as a peripheral one. The quality-assessment literature aggregates up a
weighted tree mirroring the guideline's plan structure; adopting that would
let the overall score reflect clinical priority rather than mere count.

\textbf{Cross-cohort pretraining.} Both Kineret models were trained from
scratch so that architecture and initialization would not be confounded. A
model pretrained across hospital networks and fine-tuned per site is the
natural route to something robust rather than fitted to one source, and it
would test whether the knowledge base's cohort-specific advantage survives
transfer. Building a large pretrained artifact in the spirit of HEART
\citep{huang2024heart} would open fine-tuning to other diseases and care
settings, and would be a step toward a general-purpose clinical
representation over interval data.

\textbf{Adherence as a model input at scale.} The null result here was
measured on one cohort, at one evaluation window, with adherence entering as
a mean per pattern. Richer encodings --- the per-day profile rather than its
mean, or the compliance trajectory as a sequence --- were not tried, and
whether a quality-of-care signal helps prediction is not settled by a single
aggregation.

\textbf{Mining candidate recommendations.} Because each recommendation is
declared externally as a pattern rather than implemented in analysis code,
evaluating a practice merely \emph{suspected} to matter costs no more than
evaluating the fourteen scored here. KarmaLego offers a route to finding
such candidates rather than guessing them: mine frequent interval patterns,
rank them by feature importance over the TPF representation, and translate
the survivors into TAK definitions that can then be scored and tested like
any other recommendation. That inverts the usual question --- which
practices earn a place in a guideline, and which current recommendations
track no outcome at all. Any candidate found this way is a hypothesis for
prospective evaluation, and a wide search would need more conservative
multiplicity correction than was applied here.

\textbf{Interval representation learning beyond healthcare.} INTERVenE is
built around properties of interval sequences --- that a symbol holds over
a span, that spans overlap --- rather than around anything specific to
medicine. Comparing interval representations systematically, and testing
the architecture on another interval domain such as industrial failure
prediction or process monitoring, would move the claim from architecture
engineering toward the representation question: why interval representations
work, and when they are worth their cost.

\textbf{Explainability as the primary object.} The attribution work here is
a demonstration rather than a study. A systematic comparison of explanation
quality between a semantic interval vocabulary and a non-semantic one --- the
same $\sigma$-bin contrast, judged on interpretability rather than accuracy
--- would test the central claim directly: that semantic vocabularies produce
clinically meaningful explanations at no cost in predictive performance.

\textbf{Prospective evaluation.} Everything in this thesis is retrospective.
The measurement instrument is automated, component-level, and transferable,
but whether surfacing these scores to a treating team changes practice ---
and whether changed practice changes outcomes --- is a question only a
prospective study can answer.

\section{Conclusions}
\label{sec:conclusions}

This thesis set out to measure adherence to the national inpatient diabetes
guideline automatically, to determine whether it changed after the national
implementation effort, to establish which recommendations are associated with
patient outcomes, and to shape short-horizon complication prediction as a
next-token problem over acquired medical knowledge rather than as feature-based
classification over raw measurements.

All four rest on one representation, and that is the thesis's organizing claim.
A single knowledge base and a single engine turn an OMOP record into named
abstraction intervals and, in the same pass, score how far recorded care
departed from the guideline's specifications. Prediction and measurement then
read one stream.

What the results support is narrower than that framing and more useful.
Knowledge-based abstraction is what makes a model's reasoning legible and what
makes guideline adherence measurable at national scale, while the predictive
strength belongs to the architecture --- INTERVenE-Enc leads every comparator on
both cohorts it was evaluated on. On MIMIC-IV that strength survived
substituting automatically derived intervals for the curated ones, so the
architecture is not dependent on hand-authored definitions; on the cohort the
knowledge base was written for, the same substitution cost real accuracy, and
the two results bound each other rather than settling the question. That the measured adherence signal adds no predictive
information beyond the record it is derived from is not a failure of the
measurement but what one should expect of a summary; whether following the
guideline matters clinically is a separate question, which the
propensity-weighted analysis answers affirmatively.

What stands is the clinical finding: the recommendations hardest to follow are
among those most worth following, and they share a structure --- repeated
decisions distributed across a hospitalization --- that makes them the natural
target for support at the point of care.
